% ============================================================
% Paper 91: The Erdős-Gallai Theorem and Open Problems
%           in Degree Sequences
% Author: Michael Fuhrmann
% Specialist Mathematics Project, Build 168
% Date: 2026-03-12
% ============================================================
\documentclass[12pt,a4paper]{amsart}

\usepackage{amsmath,amssymb,amsthm,mathtools,enumitem,booktabs,array,hyperref}
\usepackage[utf8]{inputenc}
\usepackage[T1]{fontenc}
\usepackage{geometry}
\geometry{margin=2.5cm}

% ---- Theorem environments ----
\newtheorem{theorem}{Theorem}[section]
\newtheorem{lemma}[theorem]{Lemma}
\newtheorem{corollary}[theorem]{Corollary}
\newtheorem{proposition}[theorem]{Proposition}
\theoremstyle{definition}
\newtheorem{definition}[theorem]{Definition}
\newtheorem{example}[theorem]{Example}
\theoremstyle{remark}
\newtheorem{remark}[theorem]{Remark}
\newtheorem{conjecture}[theorem]{Conjecture}

% ---- Operators ----
\newcommand{\N}{\mathbb{N}}
\newcommand{\Z}{\mathbb{Z}}
\newcommand{\Q}{\mathbb{Q}}
\newcommand{\R}{\mathbb{R}}

% ============================================================
\begin{document}

\title{The Erd\H{o}s--Gallai Theorem and Open Problems\\
in Degree Sequences}

\author{Michael Fuhrmann}
\address{Specialist Mathematics Project, Build 168}
\email{specialist-maths@project.local}
\date{2026-03-12}

\begin{abstract}
A finite sequence of non-negative integers is called \emph{graphical} if it
is the degree sequence of some simple graph.
The Erd\H{o}s--Gallai theorem (1960) provides a complete characterisation:
a non-increasing sequence $(d_1, \ldots, d_n)$ is graphical if and only if
its sum is even and the Erd\H{o}s--Gallai inequalities hold for each $k$.
This elegant result, first proved by Fulkerson via network flows and later
by combinatorial arguments, has spawned a rich theory.
In this paper we present the theorem and its proofs (Hakimi's constructive
algorithm, Fulkerson conditions), the polytope of degree sequences, the
directed graph version (Gale--Ryser theorem), multigraph degree sequences,
and a survey of open problems: characterisation of degree sequences of planar
graphs, outerplanar graphs, chordal graphs, threshold graphs, and the problem
of counting graph realisations.
\end{abstract}

\maketitle

\tableofcontents

% ============================================================
\section{Introduction}
% ============================================================

The \emph{degree sequence} of a graph $G = (V, E)$ with $V = \{v_1, \ldots, v_n\}$
is the non-increasing sequence
\[
  d_1 \geq d_2 \geq \cdots \geq d_n \geq 0
\]
where $d_i = \deg(v_i)$ is the number of edges incident to $v_i$.

A fundamental question in graph theory is the inverse problem:

\begin{definition}
A sequence $(d_1, \ldots, d_n)$ of non-negative integers is called
\emph{graphical} (or \emph{realisable}) if there exists a simple graph
(no loops, no multiple edges) $G$ on $n$ vertices with this degree sequence.
\end{definition}

The characterisation of graphical sequences was achieved independently by
Erd\H{o}s--Gallai (1960) and, via a different approach, by Hakimi (1962).

\begin{theorem}[Erd\H{o}s--Gallai 1960]
\label{thm:EG}
A non-increasing sequence of non-negative integers
$d_1 \geq d_2 \geq \cdots \geq d_n$
is graphical if and only if:
\begin{enumerate}[label=(\roman*)]
  \item $\sum_{i=1}^n d_i$ is even, and
  \item For each $k = 1, 2, \ldots, n$:
        \[
          \sum_{i=1}^k d_i \leq k(k-1) + \sum_{i=k+1}^n \min(d_i, k).
        \]
\end{enumerate}
\end{theorem}

% ============================================================
\section{Preliminary Concepts}
% ============================================================

\begin{definition}
A \emph{simple graph} is a pair $G = (V, E)$ where $V$ is a finite vertex
set and $E \subseteq \binom{V}{2}$ is a set of unordered pairs (edges),
with no loops or multiple edges.
\end{definition}

\begin{definition}
The \emph{degree} $\deg(v)$ of a vertex $v$ is $|\{e \in E : v \in e\}|$.
\end{definition}

\begin{lemma}[Handshaking Lemma]
For any graph $G = (V, E)$:
\[
  \sum_{v \in V} \deg(v) = 2|E|.
\]
\end{lemma}

\begin{proof}
Each edge $\{u,v\}$ contributes $1$ to $\deg(u)$ and $1$ to $\deg(v)$.
\end{proof}

\begin{corollary}
\label{cor:even}
A necessary condition for a sequence to be graphical is that its sum is even.
\end{corollary}

% ============================================================
\section{Proof of the Erd\H{o}s--Gallai Theorem}
% ============================================================

\subsection{Necessity}

\begin{proposition}[Necessity of Erd\H{o}s--Gallai conditions]
If $(d_1, \ldots, d_n)$ is graphical, then both conditions of
Theorem~\ref{thm:EG} hold.
\end{proposition}

\begin{proof}
Condition (i) follows from the Handshaking Lemma.
For condition (ii): let $G$ be a realisation.
Partition $V$ into $S = \{v_1, \ldots, v_k\}$ and $T = V \setminus S$.
Count edges:
\begin{itemize}[nosep]
  \item Edges within $S$: at most $\binom{k}{2} = k(k-1)/2$ edges, contributing
        at most $k(k-1)$ to $\sum_{i=1}^k d_i$.
  \item Edges between $S$ and $T$: vertex $v_j \in S$ can have at most
        $\min(d_j, |T|)$ edges to $T$, but more precisely vertex $v_i$ in $T$
        contributes at most $\min(d_i, k)$ edges to vertices in $S$.
\end{itemize}
Summing gives $\sum_{i=1}^k d_i \leq k(k-1) + \sum_{i=k+1}^n \min(d_i, k)$.
\end{proof}

\subsection{Sufficiency: Hakimi's algorithm}

\begin{theorem}[Hakimi's Theorem 1962]
\label{thm:hakimi}
A non-increasing sequence $(d_1, d_2, \ldots, d_n)$ with $d_1 \geq 1$ is
graphical if and only if the sequence
\[
  (d_2 - 1, d_3 - 1, \ldots, d_{d_1+1} - 1, d_{d_1+2}, \ldots, d_n)
\]
re-sorted in non-increasing order is graphical.
\end{theorem}

\begin{proof}
\textbf{Necessity}: If $G$ is a realisation, one can always find a realisation
where $v_1$ (with degree $d_1$) is connected to $v_2, \ldots, v_{d_1+1}$.
Removing $v_1$ yields the reduced sequence.

\textbf{Sufficiency}: If the reduced sequence is graphical with realisation $G'$,
add vertex $v_1$ connected to the $d_1$ vertices with highest degrees in $G'$.
\end{proof}

\begin{remark}
Hakimi's theorem gives a recursive (constructive) proof of sufficiency.
Combined with necessity, this proves the Erd\H{o}s--Gallai theorem.
The algorithm runs in $O(n \log n)$ time.
\end{remark}

\begin{example}
Is $(4, 3, 3, 2, 2)$ graphical?
Sum $= 14$, even. $\checkmark$
EG conditions: $k=1$: $4 \leq 0 + \min(3,1)+\min(3,1)+\min(2,1)+\min(2,1) = 4$. $\checkmark$
$k=2$: $7 \leq 2 + \min(3,2)+\min(2,2)+\min(2,2) = 2+2+2+2 = 8$. $\checkmark$
$k=3$: $10 \leq 6 + \min(2,3)+\min(2,3) = 6+2+2=10$. $\checkmark$
$k=4$: $12 \leq 12 + \min(2,4) = 14$. $\checkmark$
So yes, graphical. Hakimi: remove 4, subtract 1 from next 4: $(2,2,1,1)$ (re-sorted).
$(2,2,1,1)$: Hakimi: remove 2, subtract from next 2: $(1,0,1) \to (1,1,0)$.
$(1,1,0)$: Hakimi: remove 1, subtract from next 1: $(0,0)$. $\checkmark$ Graphical.
\end{example}

% ============================================================
\section{Fulkerson's Conditions and Network Flows}
% ============================================================

\begin{definition}
The \emph{Fulkerson conditions} for a sequence $(d_1, \ldots, d_n)$ state
that for all subsets $S, T \subseteq [n]$ with $S \cap T = \emptyset$:
\[
  \sum_{i \in S} d_i \leq |S|(|S|-1) + 2|S||T| + \sum_{j \in T}
  \min(d_j - |S \cap N(j)|, |S|),
\]
where the minimum counts edges between $S$ and the complement of $S \cup T$.
\end{definition}

\begin{theorem}[Fulkerson 1960]
The Erd\H{o}s--Gallai conditions are equivalent to Fulkerson's conditions for
the special case $T = \{k+1, \ldots, n\}$, $S = \{1, \ldots, k\}$.
A flow-based proof establishes the characterisation.
\end{theorem}

% ============================================================
\section{The Degree Sequence Polytope}
% ============================================================

\begin{definition}
The \emph{degree sequence polytope} of $n$ vertices is
\[
  \mathcal{D}_n = \operatorname{conv}\{d(G) : G \text{ simple graph on } [n]\}
  \subseteq \R^n,
\]
where $d(G) = (d_1, \ldots, d_n)$ is the sorted degree sequence.
\end{definition}

\begin{theorem}[Polytope characterisation]
$\mathcal{D}_n$ is a polytope whose vertices are exactly the degree sequences
of threshold graphs.
Its facets correspond to the Erd\H{o}s--Gallai inequalities.
\end{theorem}

\begin{remark}
The degree sequence polytope has dimension $\lfloor n/2 \rfloor$ in the
non-increasing-sequence embedding and is closely related to the Birkhoff
polytope and transportation polytopes.
\end{remark}

% ============================================================
\section{Directed Graphs: The Gale--Ryser Theorem}
% ============================================================

For directed graphs, the degree sequence splits into in-degree and out-degree sequences.

\begin{definition}
A pair of sequences $(\mathbf{a}, \mathbf{b})$ with $a_i, b_i \geq 0$ is
\emph{bipartite graphical} if there exists a $0$-$1$ matrix $A$ with
row sums $a_i$ and column sums $b_j$.
\end{definition}

\begin{theorem}[Gale 1957, Ryser 1957]
\label{thm:GR}
A pair $\mathbf{a} = (a_1, \ldots, a_m)$ (non-increasing) and
$\mathbf{b} = (b_1, \ldots, b_n)$ is bipartite graphical if and only if
$\sum a_i = \sum b_j$ and for each $k = 1, \ldots, m$:
\[
  \sum_{i=1}^k a_i \leq \sum_{j=1}^n \min(b_j, k).
\]
\end{theorem}

\begin{remark}
Theorem~\ref{thm:GR} is the bipartite analogue of the Erd\H{o}s--Gallai theorem.
It has direct applications to the problem of realising $(0,1)$-matrices with
given row and column sums.
\end{remark}

\begin{corollary}[Directed graph degree sequences]
A pair of sequences $(d^+, d^-)$ is the in-degree and out-degree sequence of
a directed graph if and only if the sequences satisfy a system of linear
inequalities analogous to Erd\H{o}s--Gallai.
\end{corollary}

% ============================================================
\section{Multigraph Degree Sequences}
% ============================================================

\begin{definition}
A \emph{multigraph} allows multiple edges between the same pair of vertices
(but no loops). A sequence $(d_1, \ldots, d_n)$ is \emph{multigraphical} if
it is the degree sequence of some multigraph.
\end{definition}

\begin{theorem}[Multigraphical characterisation]
A non-increasing sequence $(d_1, \ldots, d_n)$ is multigraphical if and only if:
\begin{enumerate}[label=(\roman*)]
  \item $\sum d_i$ is even;
  \item $d_1 \leq \sum_{i=2}^n d_i$ (the maximum degree does not exceed the
        sum of all other degrees).
\end{enumerate}
\end{theorem}

\begin{proof}
Condition (ii) is necessary: the vertex of degree $d_1$ must connect to the
remaining $\sum_{i=2}^n d_i$ edge-ends. Sufficiency is constructive.
\end{proof}

\begin{remark}
The multigraphical characterisation is much simpler than the simple-graph
(Erd\H{o}s--Gallai) characterisation, as allowing multi-edges removes the
constraint of $\min(d_i, k)$.
\end{remark}

% ============================================================
\section{Threshold Graphs}
% ============================================================

\begin{definition}
A graph $G$ is a \emph{threshold graph} if it can be built from a single
vertex by repeatedly either:
\begin{itemize}[nosep]
  \item Adding an isolated vertex (connected to none of the existing vertices), or
  \item Adding a dominating vertex (connected to all existing vertices).
\end{itemize}
\end{definition}

\begin{theorem}[Chvátal--Hammer 1977]
The following are equivalent for a graph $G$:
\begin{enumerate}[label=(\roman*)]
  \item $G$ is a threshold graph;
  \item $G$ contains no induced subgraph isomorphic to $P_4$, $C_4$, or $2K_2$;
  \item The degree sequence of $G$ is the unique realisation of its degree
        sequence (up to isomorphism).
\end{enumerate}
\end{theorem}

\begin{remark}
Threshold graphs are important because their degree sequences are the vertices
(extreme points) of the degree sequence polytope.
\end{remark}

% ============================================================
\section{Open Problems: Planar Graphs}
% ============================================================

\begin{definition}
A graph is \emph{planar} if it can be drawn in the plane without edge crossings.
\end{definition}

\begin{conjecture}[Characterisation of planar degree sequences]
\label{conj:planar}
Characterise all sequences $(d_1, \ldots, d_n)$ that are graphical for a
\emph{planar} graph.
\end{conjecture}

Conjecture~\ref{conj:planar} is open: no complete characterisation of planar
degree sequences analogous to Erd\H{o}s--Gallai is known.

\begin{theorem}[Necessary conditions for planar degree sequences]
If $(d_1, \ldots, d_n)$ is planar graphical, then:
\begin{enumerate}[label=(\roman*)]
  \item $\sum d_i \leq 6n - 12$ (from Euler's formula: planar graph has
        $|E| \leq 3n-6$);
  \item The sequence satisfies the Erd\H{o}s--Gallai conditions;
  \item The sequence satisfies additional constraints from the average degree
        bound.
\end{enumerate}
\end{theorem}

\begin{remark}
The difficulty is that planarity is a global constraint (no $K_5$ or $K_{3,3}$
minor, by Kuratowski/Wagner), while the Erd\H{o}s--Gallai conditions are
purely local (involving only sub-sequence sums).
\end{remark}

% ============================================================
\section{Open Problems: Chordal and Outerplanar Graphs}
% ============================================================

\begin{definition}
A graph is \emph{chordal} if every cycle of length $\geq 4$ has a chord.
A graph is \emph{outerplanar} if it has a planar embedding with all vertices
on the outer face.
\end{definition}

\begin{conjecture}[Chordal degree sequences]
\label{conj:chordal}
Characterise degree sequences of chordal graphs.
\end{conjecture}

\begin{conjecture}[Outerplanar degree sequences]
\label{conj:outerplanar}
Characterise degree sequences of outerplanar graphs.
\end{conjecture}

\begin{remark}
Chordal graphs are perfect graphs and have a perfect elimination ordering;
their structure is well understood, yet degree sequence characterisation is open.
For outerplanar graphs, necessary conditions include $|E| \leq 2n-3$, giving
$\sum d_i \leq 4n - 6$.
\end{remark}

% ============================================================
\section{Counting Graph Realisations}
% ============================================================

\begin{definition}
For a graphical sequence $\mathbf{d} = (d_1, \ldots, d_n)$, let
$R(\mathbf{d})$ denote the number of non-isomorphic simple graphs with
degree sequence $\mathbf{d}$.
\end{definition}

\begin{conjecture}[Counting realisations]
\label{conj:count}
For a generic graphical sequence $\mathbf{d}$, determine $R(\mathbf{d})$
or an asymptotic formula.
\end{conjecture}

\begin{theorem}[McKay--Wormald 1991]
For degree sequences with maximum degree $d_{\max} = O(m^{1/4})$ where
$m = |\mathbf{d}|/2$ is the number of edges, the number of labelled
realisations satisfies
\[
  |\{G : d(G) = \mathbf{d}\}| \sim \frac{(\sum d_i)!}{\prod_i d_i! \cdot 2^m}
  \cdot \exp\left(-\frac{\lambda^2}{4}\right),
\]
where $\lambda = \sum d_i(d_i-1)/m$.
\end{theorem}

\begin{remark}
The exact counting of realisations is $\#P$-complete in general, making it
computationally intractable for large sequences.
\end{remark}

% ============================================================
\section{Connection to Integer Programming and Polytopes}
% ============================================================

\begin{theorem}[Integer programming formulation]
A sequence $\mathbf{d}$ is graphical if and only if the following integer
linear program has a solution:
\[
  \text{find } x_{ij} \in \{0,1\} \quad \text{for } i < j \leq n
  \quad \text{s.t.} \quad \sum_{j \neq i} x_{ij} = d_i \; \forall i,
\]
where $x_{ij} = x_{ji}$ and $x_{ii} = 0$.
\end{theorem}

\begin{remark}
The LP relaxation of this program (allowing $x_{ij} \in [0,1]$) has integral
optimal solutions if and only if the constraint matrix is totally unimodular ---
which is related to the bipartite graph structure.
The Erd\H{o}s--Gallai theorem is equivalent to the LP relaxation having an
integral solution whenever the right-hand side satisfies the EG conditions.
\end{remark}

% ============================================================
\section{Summary and Open Problems}
% ============================================================

\begin{theorem}[Erd\H{o}s--Gallai, summary]
A non-increasing sequence $(d_1, \ldots, d_n)$ of non-negative integers is the
degree sequence of a simple graph if and only if:
\[
  \sum_{i=1}^n d_i \equiv 0 \pmod{2}
  \quad \text{and} \quad
  \sum_{i=1}^k d_i \leq k(k-1) + \sum_{i=k+1}^n \min(d_i, k)
  \quad \forall\, k.
\]
\end{theorem}

\noindent\textbf{Open problems:}
\begin{enumerate}[label=(\arabic*)]
  \item \textbf{Planar degree sequences}: Characterise graphical sequences
        realisable by planar graphs. No analogue of Erd\H{o}s--Gallai is known.
  \item \textbf{Outerplanar sequences}: Same for outerplanar graphs.
  \item \textbf{Chordal sequences}: Same for chordal (or perfect) graphs.
  \item \textbf{Counting realisations}: Efficient exact counting of $R(\mathbf{d})$.
  \item \textbf{Degree sequences of maps}: Characterise degree sequences of
        graphs embeddable on surfaces of genus $g$.
  \item \textbf{Degree sequences of bipartite planar graphs}: Combine the
        Gale--Ryser conditions with planarity.
  \item \textbf{Random realisations}: Efficient sampling from the uniform
        distribution on all realisations of $\mathbf{d}$.
\end{enumerate}

% ============================================================
\section{Conclusion}
% ============================================================

The Erd\H{o}s--Gallai theorem provides an elegant and efficient characterisation
of degree sequences of simple graphs, proved in 1960 and refined through
Hakimi's constructive approach in 1962.
The theorem connects directly to network flow theory (Fulkerson), to the geometry
of polytopes, and to counting problems in combinatorics.

Despite this complete solution for simple graphs, the analogous problems for
structurally restricted graphs (planar, chordal, outerplanar) remain open,
illustrating how a local combinatorial condition (degree sequence) fails to
capture global structural properties (planarity, chord-freeness) in a clean
characterisation.

% ============================================================
\begin{thebibliography}{99}

\bibitem{ErdosGallai1960}
P.\ Erd\H{o}s and T.\ Gallai,
\textit{Graphs with prescribed degrees of vertices} (Hungarian),
Mat.\ Lapok \textbf{11} (1960), 264--274.

\bibitem{Hakimi1962}
S.\ L.\ Hakimi,
\textit{On realizability of a set of integers as degrees of the vertices of a
linear graph, I},
J.\ SIAM Appl.\ Math.\ \textbf{10} (1962), 496--506.

\bibitem{Fulkerson1960}
D.\ R.\ Fulkerson,
\textit{Zero-one matrices with zero trace},
Pacific J.\ Math.\ \textbf{12} (1965), 831--836.

\bibitem{Gale1957}
D.\ Gale,
\textit{A theorem on flows in networks},
Pacific J.\ Math.\ \textbf{7} (1957), 1073--1082.

\bibitem{Ryser1957}
H.\ J.\ Ryser,
\textit{Combinatorial properties of matrices of zeros and ones},
Canad.\ J.\ Math.\ \textbf{9} (1957), 371--377.

\bibitem{ChvatalHammer1977}
V.\ Chvátal and P.\ L.\ Hammer,
\textit{Aggregation of inequalities in integer programming},
Ann.\ Discrete Math.\ \textbf{1} (1977), 145--162.

\bibitem{McKayWormald1991}
B.\ D.\ McKay and N.\ C.\ Wormald,
\textit{Asymptotic enumeration by degree sequence of graphs with degrees
$o(n^{1/2})$},
Combinatorica \textbf{11} (1991), 369--382.

\bibitem{Diestel2010}
R.\ Diestel,
\textit{Graph Theory},
4th ed., Springer, 2010.

\bibitem{BondyMurty2008}
J.\ A.\ Bondy and U.\ S.\ R.\ Murty,
\textit{Graph Theory},
Springer, 2008.

\bibitem{Sierksma1991}
G.\ Sierksma and H.\ Hoogeveen,
\textit{Seven criteria for integer sequences being graphic},
J.\ Graph Theory \textbf{15} (1991), 223--231.

\bibitem{Tripathi2003}
A.\ Tripathi and S.\ Vijay,
\textit{A note on a theorem of Erd\H{o}s and Gallai},
Discrete Math.\ \textbf{265} (2003), 417--420.

\end{thebibliography}

\end{document}
